\clearpage
\begingroup
\fancyhead[L]{\small THEORY OF THE MOTION OF THE MOON.}
\renewcommand{\headrulewidth}{0pt}
\par\noindent\hspace*{\parindent}We call the quantities \(n\), \(e\), \(\varepsilon\), \(g\), when multiplied by specified numbers,
quantities of one dimension. Products of two such quantities are quantities of
two dimensions, products of three are quantities of
three dimensions, and\,so\,on. In the formula for longitude, latitude, \(p\), and\,so\,forth,
we call all quantities of one dimension, of two, three, and\,so\,on,
parts of the first, second, third order, and\,so\,forth. It is immediately clear
that, to have the differential equation of the 2\textsuperscript{nd} order for \(\theta\) accurate through
quantities of the 5\textsuperscript{th} order inclusive, one must have the values of
\(\dfrac{B^3}{R^3}\), \(\dfrac{\mathrm{d}u^2}{\mathrm{d}v^2}\) accurate through
the second order, and the value of \(\theta\) through the 3\textsuperscript{rd}.
We have used only the parts whose arguments are \(\lambda\) and \(2E-\lambda\),
which are respectively of the first and second order. Consequently,
to have that differential equation accurate through the 5\textsuperscript{th} order,
we must also take account of the 5 parts just obtained by integration
that are of the 3\textsuperscript{rd} order. In this way the following additions arise
\begingroup\small
\setlength{\arraycolsep}{5pt}
\renewcommand{\arraystretch}{1.2}
\[
\begin{array}{l@{\qquad}l}
\multicolumn{1}{c}{\text{by virtue of the part}}&
\multicolumn{1}{c}{\text{[to the differential equation}}\\
\multicolumn{1}{c}{\text{[of }\theta\text{:]}}&
\multicolumn{1}{c}{\text{for }\theta\text{:]}}\\
&\multicolumn{1}{c}{\tfrac32nng\times}\\
+\tfrac58eeg\sin(2M-\lambda)&-\tfrac58ee\sin(2E-2M+\lambda)\\
-\tfrac98n\varepsilon g\sin(\lambda+a)&+\tfrac98n\varepsilon\sin(2E-\lambda-a)\\
+\tfrac98n\varepsilon g\sin(\lambda-a)&-\tfrac98n\varepsilon\sin(2E-\lambda+a)\\
+\tfrac38n\varepsilon g\sin(2E-\lambda+a)&-\tfrac38n\varepsilon\sin(\lambda-a)\\
-\tfrac78n\varepsilon g\sin(2E-\lambda-a)&+\tfrac78n\varepsilon\sin(\lambda+a).
\end{array}
\]
\endgroup
\noindent Consequently, the following is also added to the value of \(\theta\):
\[
+g\left\{\tfrac{15}{64}nee\sin(2E-2M+\lambda)\right\}
\]
\VIIInlineSourceHead{\SourceLeaf{637}}{}{THEORY OF THE MOTION OF THE MOON.}{631}
\noindent and the coefficients of
\[
\begin{array}{r|l}
\multicolumn{1}{c}{g\times\sin}&\\
\lambda+a&-\tfrac98n\varepsilon\left(1-\tfrac{23}{24}n\right)\\
\lambda-a&+\tfrac98n\varepsilon\left(1-\tfrac18n\right)\\
2E-\lambda+a&+\tfrac38n\varepsilon\left(1+\tfrac{19}{8}n\right)\\
2E-\lambda-a&-\tfrac78n\varepsilon\left(1+\tfrac{65}{56}n\right).
\end{array}
\]
\par\noindent\hspace*{\parindent}To compare this formula with [\textsc{Mayer}'s] tables,
several reductions must first be carried out. —

\par\noindent\hspace*{\parindent}I. We must replace \(E\), i.\,e. the distance of the Moon from a
fictitious Sun moving with a mean speed proportional to that of the Moon,
with \(E^{\prime}\), the distance of the Moon from the
true Sun.

\par\noindent\hspace*{\parindent}II. In place of \(\lambda\), the distance of the Moon in the ecliptic from a
fictitious node moving with a speed proportional to that of the Moon,
we must use \(\lambda^{\prime}\), the distance of the Moon along the orbit
from a uniformly moving node.

\par\noindent\hspace*{\parindent}III. Likewise, in place of \(M\) we must use \(M^{\prime}\), the distance from the uniformly advancing
apogee, and

\par\noindent\hspace*{\parindent}IV. In place of \(a\) we must use \(a^{\prime}\), the actual mean anomaly of the Sun.

\par\noindent\hspace*{\parindent}Finally, V. from \(\theta\), the tangent of the latitude, we must derive the latitude itself, \(\beta\).
For these purposes we have the following equations:
\[
\begin{aligned}
\text{I.}\quad E={}&E^{\prime}-2\varepsilon\sin a+\tfrac54\varepsilon\varepsilon\sin2a
 +2ne\sin M+\tfrac34nee\sin2M\\
&-\tfrac{11}{8}n^3\sin2E+\tfrac{15}{4}nne\sin(2E-M)
 +\tfrac14ngg\sin2\lambda-3nn\varepsilon\sin a
\end{aligned}
\]
\[
\begin{aligned}
\text{II.}\quad\lambda={}&\lambda^{\prime}
 -\left(\tfrac14gg-\tfrac18g^4\right)\sin2\lambda^{\prime}
 +\tfrac{1}{32}g^4\sin4\lambda\\
&-\zeta\bigl\{2e\sin M+\tfrac34ee\sin2M-\tfrac{11}{8}nn\sin2E\\
&\qquad+\tfrac{15}{4}ne\sin(2E-M)+\tfrac14gg\sin2\lambda-3n\varepsilon\sin a\bigr\}
\end{aligned}
\]
\[
\text{III.}\quad M=M^{\prime}+2e\eta\sin M
\]
\[
\text{IV.}\quad a=a^{\prime}+2ne\sin M
\]
\[
\text{V.}\quad\beta=\theta-\tfrac13\theta^3+\tfrac15\theta^5.
\]
\clearpage
\endgroup
